\section{数据瓶颈：从人工标注到教程驱动轨迹生成}

\subsection{为什么轨迹数据是核心稀缺资源}
讲者指出，多步交互轨迹远比单轮指令数据昂贵。每条高质量轨迹通常包含复杂状态切换、工具调用与错误恢复，纯人工采集难以支撑规模化训练。

\begin{importantbox}{数据稀缺不在 ``文本数量''，而在 ``可执行轨迹质量''}
Agent 数据的价值来自三点：任务上下文完整、动作可执行、结果可验证。缺任一项，训练增益都会显著下降。
\end{importantbox}

\subsection{教程驱动（Tutorial-driven）数据构造}
课程提出从互联网教程中抽取高价值任务模板：先做教程收集与结构化解析，再转为标准化任务描述与步骤候选，最后由 Agent 在环境中执行生成轨迹。

\begin{figure}[H]
\centering
\includegraphics[width=0.88\textwidth]{frame-05-tutorial-pipeline.jpg}
\caption{视频约 00:44:40：教程驱动数据构造流程，包含抽取、结构化、执行与筛选}
\end{figure}

\begin{knowledgebox}{教程驱动方案的优势}
\begin{itemize}
    \item 任务多样性高：覆盖真实用户问题分布；
    \item 可迁移性好：同类教程可映射为不同环境任务；
    \item 成本可控：减少纯人工逐步标注；
    \item 易扩展：可与自动评估脚本打通形成闭环。
\end{itemize}
\end{knowledgebox}

\subsection{质量控制：从 ``可运行'' 到 ``可学习''}
教程文本并不天然等于训练数据。课程强调多级过滤：任务可执行性检查、轨迹一致性检查、错误步骤剔除、低信息样本降权。高噪声数据会显著拉低训练效率。

\begin{warningbox}{不要把 ``采集成功'' 误判为 ``数据可用''}
很多轨迹虽然执行完毕，但对模型学习几乎无增益，比如重复模板路径、缺少关键决策步骤、动作与观测对齐错误。数据治理必须进入训练主流程，而非离线补丁。
\end{warningbox}

\subsection{本章小结}
课程给出的数据路线是 ``自动采集 + 结构化治理 + 脚本验证''。这条路线不追求绝对纯净，而追求在可控成本下持续提升有效样本密度。

